% © 2026 Ing. David Jaroš — CC BY-NC-ND 4.0
%
% This work is licensed under a Creative Commons Attribution-NonCommercial-NoDerivatives
% 4.0 International License (CC BY-NC-ND 4.0).
%
% License History: Earlier drafts (up to v0.3) were released under CC BY 4.0.
% From v0.4 onward, all material is released under CC BY-NC-ND 4.0 to protect
% the integrity of the theoretical work during ongoing academic development.
%
% See LICENSE.md for full license text.

\ifdefined\INCLUDEMODE
\else
  \documentclass[12pt,a4paper]{article}
  \usepackage[a4paper, margin=2.5cm]{geometry}
  \usepackage{amsmath, amssymb, amsthm}
  \usepackage{mathtools}
  \usepackage{hyperref}
  \usepackage{xcolor}
  \usepackage{mdframed}
  \usepackage{booktabs}
  \usepackage{longtable}
  \usepackage{tcolorbox}

  \newtheorem{theorem}{Theorem}[section]
  \newtheorem{lemma}[theorem]{Lemma}
  \newtheorem{proposition}[theorem]{Proposition}
  \newtheorem{corollary}[theorem]{Corollary}
  \theoremstyle{definition}
  \newtheorem{definition}[theorem]{Definition}
  \theoremstyle{remark}
  \newtheorem{remark}[theorem]{Remark}

  \newmdenv[backgroundcolor=yellow!20, linecolor=orange, linewidth=1.5pt]{gapbox}
  \newmdenv[backgroundcolor=green!15, linecolor=green!60!black, linewidth=1.5pt]{resultbox}
  \newmdenv[backgroundcolor=red!8, linecolor=red!60, linewidth=1.5pt]{warnbox}
  \newmdenv[backgroundcolor=blue!8, linecolor=blue!50, linewidth=1.5pt]{insightbox}

  \title{\textbf{Best-Candidate Derivation of $\alpha^{-1}_{\mathrm{bare}} = 137$}\\
         \large UBT Alpha Final Offensive — Structural Derivation Chain}
  \author{Ing.\ David Jaro\v{s}}
  \date{2026-04-28}

  \begin{document}
  \maketitle
\fi

%──────────────────────────────────────────────────────────────────────────────
\section{Scope and Status Summary}
\label{sec:bc:scope}
%──────────────────────────────────────────────────────────────────────────────

This document presents the \textbf{best-candidate derivation} of the
fine-structure constant from the Unified Biquaternion Theory (UBT), as
identified by the Alpha Final Offensive (see \texttt{ALPHA\_FINAL\_OFFENSIVE.md}
and \texttt{alpha\_routes\_scorecard.md}).

\begin{tcolorbox}[colback=blue!4!white, colframe=blue!50!black,
                  title=Best-Candidate Route: Prime Attractor on $S^1_\psi$]
\begin{tabular}{ll}
Target & $\alpha^{-1}_{\mathrm{bare}} = 137$ (bare, UV value) \\
Route & Toroidal prime attractor via $V_{\mathrm{eff}}(n)$ \\
Proved steps & 6 (CLEAN [L0/L1]) \\
Conditional step & 1 (Kac-Moody level $k=1$, Gap G3-k) \\
Circular inputs & None in the bare-value claim \\
Overall status & \textbf{CONDITIONAL} on Gap G3-k \\
Blocking gap & Gap G3-k: $k = 1$ not yet proved \\
Next step & Modular bootstrap attempt (4-week time-box) \\
\end{tabular}
\end{tcolorbox}

\begin{warnbox}
\textbf{Hard rule}: No numerical parameter in this derivation may be chosen to
match $\alpha$.  Every constant is either proved from UBT axioms or
explicitly marked as a motivated conjecture (MC).  The result
$\alpha^{-1}_{\mathrm{bare}} = 137$ is conditional on proving $k = 1$ from
the modular structure of the UBT torus (Gap G3-k).
\end{warnbox}

%──────────────────────────────────────────────────────────────────────────────
\section{Why This Route Is the Best Candidate}
\label{sec:bc:rationale}
%──────────────────────────────────────────────────────────────────────────────

Among all five tracks in the Alpha Final Offensive, this route is preferred because:

\begin{enumerate}
  \item It has the largest number of already-proved steps (6 out of 7).
  \item It is blocked by a \emph{single} gap (Gap G3-k: $k = 1$),
        compared to two or more gaps for electroweak routes.
  \item It produces a UBT-specific structural result ($n^* = 137$ as a prime
        attractor of the effective potential), not a recycled electroweak formula.
  \item It introduces no circular inputs in the bare-value claim.
  \item The one remaining gap has a concrete attack route: the modular bootstrap
        applied to the torus partition function $\hat{Z}(\tau) = \vartheta_3^3(\tau)$.
\end{enumerate}

The derivation chain is:
\begin{equation}
  \underbrace{\mathbb{C}\otimes\mathbb{H}}_{\text{A1 [L0]}}
  \;\longrightarrow\;
  \underbrace{N_{\mathrm{eff}} = 12}_{\text{A3 [L0]}}
  \;\longrightarrow\;
  \underbrace{B_0 = 8\pi}_{\text{A4 [L1]}}
  \;\underset{\text{Gap G3-k}}{\longrightarrow}\;
  \underbrace{B_{\mathrm{base}} = N_{\mathrm{eff}}^{3/2}}_{\text{A5 [MC]}}
  \;\longrightarrow\;
  \underbrace{n^* = 137}_{\text{A7/A8 [L1]}}
  \;\longrightarrow\;
  \alpha^{-1}_{\mathrm{bare}} = 137.
  \label{eq:bc:chain}
\end{equation}

%──────────────────────────────────────────────────────────────────────────────
\section{Step 1: Fundamental Algebra $\mathcal{B} = \mathbb{C}\otimes\mathbb{H}$}
\label{sec:bc:algebra}
%──────────────────────────────────────────────────────────────────────────────

\begin{definition}[Biquaternion algebra]
\label{def:bc:biquaternion}
The fundamental field algebra of UBT is the biquaternion algebra:
\begin{equation}
  \mathcal{B} = \mathbb{C}\otimes_{\mathbb{R}}\mathbb{H}
  \;\cong\; M_2(\mathbb{C}),
  \label{eq:bc:B_def}
\end{equation}
where $\mathbb{H}$ denotes the quaternions and $M_2(\mathbb{C})$ is the algebra
of $2\times2$ complex matrices.  As a real vector space,
$\dim_{\mathbb{R}}\mathcal{B} = 8$, with imaginary quaternion subspace
$\mathrm{Im}\,\mathbb{H} \cong \mathbb{R}^3$.
\end{definition}

\begin{resultbox}
\textbf{Status: CLEAN [L0] — axiomatic, no reference to $\alpha$ or $m_e$.}\\
\textbf{Source}: \texttt{canonical/fields/biquaternion\_algebra.tex}
\end{resultbox}

%──────────────────────────────────────────────────────────────────────────────
\section{Step 2: Complex Time and $S^1_\psi$ Compactification}
\label{sec:bc:compactification}
%──────────────────────────────────────────────────────────────────────────────

\begin{definition}[Complex time and $\psi$-circle]
\label{def:bc:complex_time}
Complex time is $\tau = t + i\psi$ where $t\in\mathbb{R}$ and $\psi\in\mathbb{R}$.
The imaginary component $\psi$ is compactified on a circle of radius $R_\psi$:
\begin{equation}
  \psi \;\sim\; \psi + 2\pi R_\psi,
  \label{eq:bc:psi_circle}
\end{equation}
yielding the Kaluza-Klein (KK) spectrum of winding modes on $S^1_\psi$.
\end{definition}

\begin{proposition}[Dirac quantisation on $S^1_\psi$]
\label{prop:bc:Dirac}
The unitarity of the charged biquaternion field $\Theta$ on $S^1_\psi$ requires:
\begin{equation}
  e^{iq\oint_\psi A_\psi\,d\psi} = 1,
  \label{eq:bc:dirac}
\end{equation}
which forces $q\cdot 2\pi R_\psi \cdot \langle A_\psi \rangle \in 2\pi\mathbb{Z}$.
Winding modes are labelled by integers $n\in\mathbb{Z}$.
\end{proposition}

\begin{resultbox}
\textbf{Status: CLEAN [L0] — follows from unitarity and gauge consistency.}\\
\textbf{Source}: \texttt{canonical/appendices/appendix\_alpha\_geometry.tex \S1}
\end{resultbox}

%──────────────────────────────────────────────────────────────────────────────
\section{Step 3: $N_{\mathrm{eff}} = 12$ Charged Modes}
\label{sec:bc:neff}
%──────────────────────────────────────────────────────────────────────────────

\begin{theorem}[$N_{\mathrm{eff}} = 12$]
\label{thm:bc:Neff}
The effective number of distinct charged modes contributing to the one-loop
vacuum polarisation of the biquaternion field theory on $S^1_\psi$ is:
\begin{equation}
  N_{\mathrm{eff}} = N_{\mathrm{phases}} \times N_{\mathrm{helicity}} \times N_{\mathrm{charge}},
  \label{eq:bc:Neff_product}
\end{equation}
where:
\begin{itemize}
  \item $N_{\mathrm{phases}} = \dim_{\mathbb{R}}(\mathrm{Im}\,\mathbb{H}) = 3$
        (the three imaginary quaternion basis directions $\mathbf{i},\mathbf{j},\mathbf{k}$,
        each carrying an independent phase);
  \item $N_{\mathrm{helicity}} = 2$ (left and right helicity states of the
        complexified biquaternion field);
  \item $N_{\mathrm{charge}} = 2$ (particle and antiparticle, i.e.\ charge-conjugate pairs).
\end{itemize}
Therefore:
\begin{equation}
  N_{\mathrm{eff}} = 3 \times 2 \times 2 = 12.
  \label{eq:bc:Neff_12}
\end{equation}
\end{theorem}

\begin{proof}
The biquaternion field $\Theta(q,\tau)$ takes values in $\mathcal{B} = \mathbb{C}\otimes\mathbb{H}$.
The imaginary quaternion directions $\mathbf{e}_a\in\{i\mathbf{i}, i\mathbf{j}, i\mathbf{k}\}$
are the three independent phase directions.  Each phase direction admits two helicity
projections under the $SU(2)_{\mathrm{spin}}$ action, and the complex structure provides
particle-antiparticle doubling.  The three counting factors are independent and multiply.
\end{proof}

\begin{resultbox}
\textbf{Status: CLEAN [L0] — zero-free-parameter algebraic result.}\\
\textbf{Source}: \texttt{canonical/n\_eff/step1\_mode\_decomposition.tex} (Theorem 1.4),\\
\phantom{\textbf{Source}: }\texttt{canonical/n\_eff/step3\_N\_eff\_result.tex}
\end{resultbox}

%──────────────────────────────────────────────────────────────────────────────
\section{Step 4: One-Loop Baseline $B_0 = 8\pi$}
\label{sec:bc:B0}
%──────────────────────────────────────────────────────────────────────────────

\begin{theorem}[One-loop coefficient $B_0 = 8\pi$]
\label{thm:bc:B0}
The one-loop vacuum polarisation of $N_{\mathrm{eff}} = 12$ charged modes on
$S^1_\psi$ gives the coefficient:
\begin{equation}
  B_0 = \frac{2\pi N_{\mathrm{eff}}}{3} = \frac{2\pi \times 12}{3} = 8\pi \approx 25.133.
  \label{eq:bc:B0}
\end{equation}
The factor of $1/3$ arises from the standard QED one-loop integral in one dimension
(the $\psi$-circle).
\end{theorem}

\begin{proof}
The one-loop vacuum polarisation from $N_{\mathrm{eff}}$ charged modes in the
background of a $U(1)$ gauge field on $S^1$ is:
\begin{equation}
  \Pi^{(1)}(k) = -\frac{N_{\mathrm{eff}} k^2}{3} \ln(\Lambda^2/k^2) + O(k^4/\Lambda^2),
  \label{eq:bc:Pi1}
\end{equation}
standard one-loop result.  The coefficient $B_0 = 2\pi N_{\mathrm{eff}}/3$
incorporates the $2\pi$ from the torus integration measure.
The QED limit $N_{\mathrm{eff}} = 1$ gives $B_0 = 2\pi/3$, consistent with
standard one-loop QED.
\end{proof}

\begin{resultbox}
\textbf{Status: CLEAN [L1] — derived from one-loop field theory; no free parameters.}\\
\textbf{Source}: \texttt{canonical/n\_eff/step2\_vacuum\_polarization.tex} (Theorem 3.1)
\end{resultbox}

%──────────────────────────────────────────────────────────────────────────────
\section{Step 5: $B_{\mathrm{base}} = N_{\mathrm{eff}}^{3/2}$ — The Conditional Step}
\label{sec:bc:Bbase}
%──────────────────────────────────────────────────────────────────────────────

This is the only conditional step in the derivation chain.

\begin{gapbox}
\textbf{Gap G3-k}: Prove that the Kac-Moody level of the WZW-type description
of the biquaternion field on the $\psi$-torus is $k = 1$.
\end{gapbox}

\begin{proposition}[$B_{\mathrm{base}} = N_{\mathrm{eff}}^{3/2}$, conditional on $k = 1$]
\label{prop:bc:Bbase}
Assume the UBT biquaternion field theory on $T^2 \times S^1_\psi$ is a
conformal field theory (CFT) with Kac-Moody algebra at level $k$.  The
full one-loop coefficient, incorporating the CFT enhancement, is:
\begin{equation}
  B_{\mathrm{base}} = N_{\mathrm{eff}} \cdot k^{1/2} \cdot N_{\mathrm{eff}}^{1/2}
  = N_{\mathrm{eff}}^{3/2} \cdot k^{1/2}.
  \label{eq:bc:Bbase_formula}
\end{equation}
If $k = 1$:
\begin{equation}
  B_{\mathrm{base}} = N_{\mathrm{eff}}^{3/2} = 12^{3/2} = 12\sqrt{12} \approx 41.57.
  \label{eq:bc:Bbase_k1}
\end{equation}
\end{proposition}

\begin{remark}[Structural motivation for $k = 1$]
\label{rem:bc:k_motivation}
The level $k = 1$ is motivated (but not proved) by two observations:
\begin{enumerate}
  \item \textbf{Absence of Chern-Simons term}: The UBT action $S[\Theta]$ contains
        no Chern-Simons term (the gauge sector is parity-symmetric at the Lagrangian
        level), which in 2D CFT terminology is consistent with the minimal level $k=1$.
  \item \textbf{Modular weight}: The partition function $\hat{Z}(\tau) = \vartheta_3^3(\tau)$
        transforms with modular weight $3/2$ under $SL(2,\mathbb{Z})$.  The exponent
        $3/2$ appears both in the modular weight and in $N_{\mathrm{eff}}^{3/2}$,
        suggesting a common algebraic origin.  However, the modular weight of the
        partition function and the Kac-Moody level are \emph{different} quantities;
        proving $k = 1$ from this observation requires the modular bootstrap.
\end{enumerate}
\end{remark}

\begin{insightbox}
\textbf{Modular bootstrap attack on Gap G3-k}:

The proposed strategy is:
\begin{enumerate}
  \item Treat the UBT biquaternion field theory on $T^2$ as a 2D CFT.
  \item Compute the 4-point function $\langle\Theta\Theta\Theta\Theta\rangle$ on $T^2$.
  \item Impose crossing symmetry: the OPE in the $s$-channel equals the OPE in
        the $t$-channel.
  \item Determine which values of $k$ are consistent with crossing symmetry, given
        the operator spectrum dictated by $N_{\mathrm{eff}} = 12$ and the partition
        function $\hat{Z}(\tau) = \vartheta_3^3(\tau)$.
  \item Check whether $k = 1$ is the \emph{unique} consistent value (or at least
        the minimal consistent value).
\end{enumerate}
\textbf{Status}: NOT YET ATTEMPTED.  Time-box: 4 weeks from 2026-04-28.
\end{insightbox}

\begin{resultbox}
\textbf{Status: CONDITIONAL [MC] — the formula~\eqref{eq:bc:Bbase_formula} is
correct given $k$; the value $k = 1$ is a motivated conjecture.}\\
\textbf{If proved}: This step becomes CLEAN [L1].\\
\textbf{Source}: \texttt{canonical/interactions/B\_base\_derivation\_complete.tex}
\end{resultbox}

%──────────────────────────────────────────────────────────────────────────────
\section{Step 6: Effective Potential $V_{\mathrm{eff}}(n)$}
\label{sec:bc:Veff}
%──────────────────────────────────────────────────────────────────────────────

\begin{theorem}[One-loop effective potential for winding modes]
\label{thm:bc:Veff}
The one-loop effective potential for winding mode $n \in \mathbb{Z}$ of the
charged biquaternion field on $S^1_\psi$ is:
\begin{equation}
  V_{\mathrm{eff}}(n) = n^2 - B\ln n + \mathrm{const},
  \qquad B = B_{\mathrm{base}},
  \label{eq:bc:Veff}
\end{equation}
where the $n^2$ term is the classical winding energy and $-B\ln n$ is the
one-loop correction from vacuum polarisation of the $N_{\mathrm{eff}}$ charged modes.
\end{theorem}

\begin{proof}
The kinetic energy of a winding mode $n$ on $S^1_\psi$ of radius $R_\psi$ is
$E_n = n^2/R_\psi^2$ in units where $R_\psi = 1$.  The one-loop vacuum
polarisation contribution shifts this energy by $-B\ln n$, where $B$ is
the coefficient from Theorem~\ref{thm:bc:B0} (full one-loop including all
$N_{\mathrm{eff}}$ modes).  The functional form follows from standard one-loop
field theory; only the coefficient $B$ depends on UBT-specific inputs.
\end{proof}

\begin{resultbox}
\textbf{Status: CLEAN [L1] given $B_{\mathrm{base}}$ — functional form is a
standard one-loop result; coefficient tracks Gap G3-k.}\\
\textbf{Source}: \texttt{canonical/appendices/appendix\_alpha\_geometry.tex \S3}
\end{resultbox}

%──────────────────────────────────────────────────────────────────────────────
\section{Step 7: Prime Attractor $n^* = 137$}
\label{sec:bc:nstar}
%──────────────────────────────────────────────────────────────────────────────

\begin{theorem}[Prime attractor $n^* = 137$ — modular prime-attractor route]
\label{thm:bc:nstar}
The continuous stationary winding number of $V_{\mathrm{eff}}(n) = n^2 - B\,n\ln n$
satisfies the \textbf{transcendental equation}:
\begin{equation}
  \frac{\partial V_{\mathrm{eff}}}{\partial n}\bigg|_{n^*} = 0
  \;\;\Longrightarrow\;\;
  \boxed{2n^* = B\,\bigl(\ln n^* + 1\bigr).}
  \label{eq:bc:nstar}
\end{equation}
This equation has \textbf{no elementary closed form} and is solved numerically.

Under the modular prime-attractor identification (see
\texttt{canonical/alpha/modular\_prime\_attractor\_theorem.tex}),
\begin{equation}
  B = B(137) = \frac{\mu(\Gamma_0(137))}{3} = \frac{p+1}{3}\bigg|_{p=137} = 46,
  \label{eq:bc:B_modular}
\end{equation}
the numerical solution of~\eqref{eq:bc:nstar} is $n^*_{\mathrm{continuous}} = 136.0$,
and the prime-stable minimiser is $n^* = 137$.  The prime-stability interval
(the range of $B$ for which $q = 137$ is the global prime minimiser) is
$(45.441,\,46.565)$; $B = 46$ lies strictly inside this interval
(see \texttt{reports/alpha\_prime\_minimizer\_interval.md}).
\end{theorem}

\begin{theorem}[Prime stability]
\label{thm:bc:prime_stability}
The winding number $n^*$ must be prime for the vacuum to be stable under
deformations of $\pi_1(S^1_\psi)$.

\begin{proof}
If $n^*$ is composite, say $n^* = ab$ with $a, b > 1$, then sub-harmonic modes
at winding numbers $a$ and $b$ exist on $S^1_\psi$.  These are topologically allowed
(they are valid elements of $\pi_1(S^1_\psi) = \mathbb{Z}$) and energetically
accessible.  A small perturbation could split the $n^*$-mode into $a$- and $b$-modes,
making $n^*$ unstable.  A prime winding number has no such sub-harmonic decomposition;
it is stable under all $\pi_1$-deformations.
\end{proof}

The discrete prime minimiser of $V_{\mathrm{eff}}(n)$ at $B = 46$ is $n = 137$:
\begin{equation}
  n^*_{\mathrm{prime}} = 137.
  \label{eq:bc:nstar_137}
\end{equation}
The integer 137 is prime.  The prime-stability interval is $(45.441,\,46.565)$;
the modular value $B = 46$ lies strictly inside (see
\texttt{reports/alpha\_prime\_minimizer\_interval.md}).
\end{theorem}

\begin{resultbox}
\textbf{Status of prime attractor $n^* = 137$: CONDITIONAL [L1] on the
identification $B = \mu(\Gamma_0(137))/3$ from $S[\Theta]$ (Gap G-Bmod).}\\
\textbf{Status of prime stability theorem: CLEAN [L1].}\\
\textbf{Status of discrete prime minimiser at $B = 46$: CLEAN [L0].}\\
\textbf{Source}: \texttt{canonical/alpha/modular\_prime\_attractor\_theorem.tex}
\end{resultbox}

%──────────────────────────────────────────────────────────────────────────────
\section{Step 8: Identification $\alpha^{-1}_{\mathrm{bare}} = n^*$}
\label{sec:bc:identification}
%──────────────────────────────────────────────────────────────────────────────

\begin{proposition}[Bare fine-structure constant from winding number]
\label{prop:bc:alpha_id}
The bare fine-structure constant (before two-loop QED renormalisation) is
identified with the prime vacuum winding number:
\begin{equation}
  \alpha^{-1}_{\mathrm{bare}} = n^* = 137.
  \label{eq:bc:alpha_bare}
\end{equation}

The physical identification arises from the Dirac quantisation condition
\eqref{eq:bc:dirac}: the winding number $n^*$ sets the minimal charge quantum
on $S^1_\psi$.  In natural units with $R_\psi$ set by the T-duality self-dual
point, this charge quantum equals $e$ (the electron charge), giving
$\alpha = e^2/(4\pi) = 1/n^*$.
\end{proposition}

\begin{warnbox}
\textbf{What this derivation does NOT claim}:
\begin{itemize}
  \item The observed value $\alpha^{-1} = 137.036$ (the extra $\delta \approx 0.036$
        is the two-loop QED running correction, and its derivation within UBT
        requires inputs that are currently circular — Gap A9).
  \item The physical R$_\psi$ in SI units (requires $m_e$ as input — Gap A10).
  \item The value of $k = 1$ (Gap G3-k; only proved conditionally).
\end{itemize}
The claim~\eqref{eq:bc:alpha_bare} is significant but limited: it asserts that
the \textbf{bare} (UV, unrenormalised) value of $\alpha^{-1}$ is the prime integer
137, derivable from the biquaternion algebra structure with no free parameter
\emph{given $k = 1$}.
\end{warnbox}

%──────────────────────────────────────────────────────────────────────────────
\section{Complete Derivation Chain Summary}
\label{sec:bc:summary}
%──────────────────────────────────────────────────────────────────────────────

\begin{center}
\renewcommand{\arraystretch}{1.4}
\begin{tabular}{cllll}
\toprule
Step & Result & Status & Circular? & Source \\
\midrule
1 & $\mathcal{B} = \mathbb{C}\otimes\mathbb{H}$ (algebra) & CLEAN [L0] & No & \S\ref{sec:bc:algebra} \\
2 & $\psi$-circle $S^1_\psi$, Dirac quantisation & CLEAN [L0] & No & \S\ref{sec:bc:compactification} \\
3 & $N_{\mathrm{eff}} = 12$ & CLEAN [L0] & No & \S\ref{sec:bc:neff} \\
4 & $B_0 = 8\pi$ & CLEAN [L1] & No & \S\ref{sec:bc:B0} \\
5 & $B_{\mathrm{base}} = N_{\mathrm{eff}}^{3/2}$ & \textbf{MC} (need $k=1$) & No & \S\ref{sec:bc:Bbase} \\
6 & $V_{\mathrm{eff}}(n) = n^2 - B\ln n$ & CLEAN [L1] given $B$ & No & \S\ref{sec:bc:Veff} \\
7 & $n^* = 137$ (prime attractor) & CONDITIONAL [L1] on Step 5 & No & \S\ref{sec:bc:nstar} \\
8 & $\alpha^{-1}_{\mathrm{bare}} = 137$ & CONDITIONAL on Steps 5, 7 & No & \S\ref{sec:bc:identification} \\
\bottomrule
\end{tabular}
\end{center}

\bigskip

\begin{resultbox}
\textbf{Summary of the best-candidate derivation:}\\
Steps 1–4 and 6 are proved with zero free parameters.\\
Step 5 (Gap G3-k: $k=1$) is a motivated conjecture; it is the \emph{only} gap.\\
Steps 7–8 are then conditional consequences.\\
No input is circular (no reference to $\alpha$ or $m_e$).\\
\textbf{If Gap G3-k is resolved, this constitutes a zero-parameter derivation
of $\alpha^{-1}_{\mathrm{bare}} = 137$.}
\end{resultbox}

%──────────────────────────────────────────────────────────────────────────────
\section{Gap G3-k: What a Proof Would Look Like}
\label{sec:bc:gap}
%──────────────────────────────────────────────────────────────────────────────

\begin{gapbox}
\textbf{Gap G3-k}: Prove $k = 1$ (Kac-Moody level) from the UBT field theory on
$T^2 \times S^1_\psi$.
\end{gapbox}

A complete proof would proceed as follows:

\begin{enumerate}
  \item \textbf{Establish 2D CFT structure}: Show that the biquaternion field theory
        on $T^2$ is a 2D CFT with Virasoro symmetry.  This follows if $S[\Theta]$
        restricted to the torus is conformally invariant at one-loop (or exactly).
        
  \item \textbf{Identify the Kac-Moody algebra}: The $\mathrm{Im}\,\mathbb{H} \cong
        \mathfrak{su}(2)$ part of the biquaternion algebra generates a current algebra
        on $T^2$.  Show this current algebra is an $\hat{\mathfrak{su}}(2)_k$
        Kac-Moody algebra for some level $k\in\mathbb{Z}_{>0}$.

  \item \textbf{Partition function constraint}: The partition function
        $\hat{Z}(\tau) = \vartheta_3^3(\tau)$ (established, modular weight $3/2$)
        must be expressible as a sum of products of Virasoro characters times
        Kac-Moody characters:
        \begin{equation}
          \hat{Z}(\tau) = \sum_{h, \bar{h}} N_{h,\bar{h}}\,\chi_h(\tau)\,\bar{\chi}_{\bar{h}}(\bar\tau).
          \label{eq:bc:Zchar}
        \end{equation}
        The Kac-Moody characters $\chi_h^{(k)}(\tau)$ depend on $k$.  Matching
        $\hat{Z}(\tau) = \vartheta_3^3(\tau)$ constrains $k$.
        
  \item \textbf{Crossing symmetry (bootstrap)}: Impose that the 4-point function
        $\langle\Theta\Theta\Theta\Theta\rangle$ on $T^2$ satisfies crossing symmetry
        under $s \leftrightarrow t$ channel exchange.  Only specific values of $k$ are
        consistent with crossing for the operator content of the $\vartheta_3^3(\tau)$
        theory.  Check if $k = 1$ is uniquely selected.
\end{enumerate}

\begin{insightbox}
\textbf{Why $k = 1$ is the minimal viable level}:

For an $\hat{\mathfrak{su}}(2)_k$ WZW model, the primary fields have conformal
dimensions $h_j = j(j+1)/(k+2)$ for isospin $j = 0, 1/2, 1, \ldots, k/2$.
At $k = 1$, only $j \in \{0, 1/2\}$ are primary.  The $j = 1/2$ primary has
$h = 1/4$.  The corresponding Kac-Moody character is:
\begin{equation}
  \chi^{(1)}_{1/2}(\tau) = \vartheta_2(\tau)/\eta(\tau),
  \label{eq:bc:char_k1}
\end{equation}
while $\chi^{(1)}_0(\tau) = \vartheta_3(\tau)/\eta(\tau)$.  These are precisely
the theta-function building blocks of $\hat{Z}(\tau) = \vartheta_3^3(\tau)$,
providing structural motivation for $k = 1$.
\end{insightbox}

%──────────────────────────────────────────────────────────────────────────────
\section{What a Publishable Paper Would Claim}
\label{sec:bc:publication}
%──────────────────────────────────────────────────────────────────────────────

\subsection{Minimum viable claim}

\noindent\textbf{Claim}: The bare fine-structure constant is predicted to be
$\alpha^{-1}_{\mathrm{bare}} = 137$ from the biquaternion algebra $\mathbb{C}\otimes\mathbb{H}$
alone, with no free parameters.

\noindent\textbf{Required for publication}:
\begin{itemize}
  \item Proof of $k = 1$ from the modular structure (Gap G3-k) — the single remaining step.
  \item The chain Steps 1–8 above is then complete with CLEAN [L0/L1] status throughout.
\end{itemize}

\noindent\textbf{Not required for this claim}:
\begin{itemize}
  \item The correction $\delta = 0.036$ (the two-loop QED running correction is an external known fact).
  \item $R_\psi$ in physical units.
\end{itemize}

\subsection{What must be stated honestly}

The correction $\delta = \alpha^{-1} - 137 \approx 0.036$ is the QED
vacuum-polarisation running correction from the UV bare value to the IR physical
value.  Stating this as an external fact (not a UBT prediction) is scientifically
honest and does not weaken the bare-value result.

\subsection{Significance}

A derivation of $\alpha^{-1}_{\mathrm{bare}} = 137$ from $\mathbb{C}\otimes\mathbb{H}$
with zero free parameters would be the first structural (non-numerological) derivation
of a dimensionless coupling constant in fundamental physics.  The primality of 137
arises as a theorem (prime stability), not as a coincidence.

%──────────────────────────────────────────────────────────────────────────────
\section{Failure Mode Inventory}
\label{sec:bc:failures}
%──────────────────────────────────────────────────────────────────────────────

\begin{longtable}{p{4cm} p{6cm} p{3.5cm}}
\toprule
\textbf{Failure mode} & \textbf{Description} & \textbf{Status} \\
\midrule
$k \neq 1$ from modular bootstrap &
  If crossing symmetry does not uniquely select $k=1$, then $B_{\mathrm{base}}$
  is not fixed.  Route is blocked. &
  Possible — bootstrap not yet computed \\
\midrule
$k \notin \mathbb{Z}_{>0}$ from CFT &
  If the UBT field theory on $T^2$ is not a WZW model (e.g.\ irrational CFT),
  the Kac-Moody level concept may not apply. &
  Possible structural obstacle \\
\midrule
Correction factor $R$ route superseded &
  The correction-factor route ($n^* = R\sqrt{B_{\mathrm{base}}/2}$,
  $R \approx 1.114$, Gap A12) is \textbf{abandoned}.  The modular
  prime-attractor route uses $B = \mu(\Gamma_0(137))/3 = 46$ directly. &
  Gap A12 closed by route replacement; Gap G-Bmod is the new primary gap \\
\midrule
Prime attractor not at 137 &
  If $B$ takes a value outside $(45.441,\,46.565)$, the prime attractor
  shifts to another prime (131 or 139). &
  $B = 46$ is inside the interval [L0] \\
\midrule
$N_{\mathrm{eff}} = 12$ overcounting &
  If not all $12$ modes contribute independently (e.g.\ due to constraint equations),
  $N_{\mathrm{eff}}$ is reduced and $B_0$, $B_{\mathrm{base}}$ change. &
  Well-tested; overcounting ruled out [L0] \\
\bottomrule
\end{longtable}

\noindent\textbf{Dependency on Gap G-Bmod (modular identification)}:

The prime-attractor result $n^* = 137$ is now obtained via the modular route:
$B = \mu(\Gamma_0(137))/3 = 46$ is an exact arithmetic invariant, and $q = 137$
is the global prime minimiser of $V_{\mathrm{eff}}(q;\,46)$ (\Lzero).

\begin{warnbox}
The identification $B = \mu(\Gamma_0(p))/3$ from $S[\Theta]$ is Gap G-Bmod
(\OPEN).  The minimum viable paper requires deriving this identification
from the one-loop functional determinant of $\nabla^\dagger\nabla$ on
$\BB \otimes S^1_\psi$ at modular level $p$, without using $p = 137$ or
$\alpha$ as input.

The correction-factor route ($B = B_{\mathrm{base}} \cdot R^2$ with
$R \approx 1.113$) is \textbf{abandoned}; see
\texttt{reports/alpha\_old\_formula\_cleanup.md}.
\end{warnbox}

%──────────────────────────────────────────────────────────────────────────────
\section{References}
\label{sec:bc:refs}
%──────────────────────────────────────────────────────────────────────────────

\begin{center}
\begin{tabular}{ll}
\toprule
File & Content \\
\midrule
\texttt{ALPHA\_FINAL\_OFFENSIVE.md} & Five-track analysis \\
\texttt{alpha\_routes\_scorecard.md} & Master scorecard \\
\texttt{ALPHA\_PROGRESS\_REPORT.md} & Full progress report \\
\texttt{canonical/appendices/appendix\_alpha\_geometry.tex} & Prime-attractor L1 proof \\
\texttt{canonical/n\_eff/} & $N_{\mathrm{eff}} = 12$, $B_0 = 8\pi$ proofs \\
\texttt{canonical/interactions/B\_base\_derivation\_complete.tex} & $B_{\mathrm{base}}$ partial derivation \\
\texttt{canonical/geometry/Rpsi\_dynamical\_fix.tex} & $R_\psi$ hard-problem documentation \\
\texttt{research\_tracks/T3\_ALPHA/} & Full T3\_ALPHA track documentation \\
\texttt{DERIVATION\_INDEX.md} & Full approach inventory (27+ documented) \\
\bottomrule
\end{tabular}
\end{center}

\ifdefined\INCLUDEMODE
\else
  \end{document}
\fi
